\begin{abstract}
    Public health agencies worldwide detect geographic disease clusters
    by computing the Kulldorff spatial scan statistic over case counts
    reported by clinics, laboratories, or sentinel sites. Publishing the
    results creates disclosure risk: a detected cluster containing only
    two or three reporting nodes effectively identifies those nodes by
    name. Despite decades of operational deployment through the SaTScan
    software, no formal privacy mechanism exists for the scan statistic.

    We provide the first differentially private release mechanism for
    the Kulldorff log-likelihood ratio. We derive a closed-form upper
    bound on the sensitivity of the scan statistic as a function of
    cluster density~$k_w$ (Theorem~1), then use this bound to construct
    an exponential mechanism whose normalized utility has sensitivity~$1$
    by construction (Theorem~2). For the prospective surveillance
    setting, where results are released weekly, we apply
    zero-concentrated differential privacy composition to show that a
    fixed annual privacy budget supports approximately
    $T_{\max} \approx [Z]$ weekly releases under zCDP, compared to
    roughly $10$ under na\"ive composition (Theorem~3).

    Experiments on synthetic outbreak data over a realistic
    province-level reporting network show that at $\varepsilon = 1$,
    the mechanism detects $[X]\%$ of injected clusters while the
    median utility gap remains below $[Y]\%$ of the non-private optimum.
    These results demonstrate that formal privacy protection is
    compatible with operationally useful disease surveillance.
\end{abstract}